% =====================================================================
% Perceiver & Perceiver IO -- Dispensa da presentazione
% Struttura: Note presentazione.pptx / presentation_v2_eng.pdf
% Contenuti: slide V2, paper Perceiver, paper Perceiver IO, appunti locali
% Compilare con: latexmk -pdf perceiver_note_presentazione.tex
% =====================================================================
\documentclass[11pt,a4paper]{article}

\usepackage[T1]{fontenc}
\usepackage[utf8]{inputenc}
\usepackage[italian]{babel}
\usepackage{lmodern}
\IfFileExists{microtype.sty}{\usepackage{microtype}}{}
\usepackage[a4paper,margin=2.25cm,headheight=15pt]{geometry}
\usepackage{amsmath,amssymb,mathtools}
\usepackage{booktabs,array,tabularx,longtable}
\usepackage{enumitem}
\usepackage[dvipsnames]{xcolor}
\usepackage[most]{tcolorbox}
\usepackage{tikz}
\usetikzlibrary{arrows.meta,positioning,calc,fit,shapes.geometric}
\usepackage{fancyhdr}
\usepackage{hyperref}

\definecolor{mainblue}{HTML}{123C69}
\definecolor{deepblue}{HTML}{0B2545}
\definecolor{softblue}{HTML}{E8F1FA}
\definecolor{softgreen}{HTML}{E9F7EF}
\definecolor{softamber}{HTML}{FFF6DE}
\definecolor{softred}{HTML}{FCEBEA}
\definecolor{accentgreen}{HTML}{1B7F5A}
\definecolor{accentamber}{HTML}{B7791F}
\definecolor{accentred}{HTML}{B3392F}

\hypersetup{
  colorlinks=true,
  linkcolor=mainblue,
  urlcolor=mainblue,
  citecolor=mainblue,
  pdftitle={Perceiver e Perceiver IO -- Dispensa da presentazione},
  pdfauthor={Francesco Gigli e Corso Vignoli}
}

\pagestyle{fancy}
\fancyhf{}
\fancyhead[L]{\small Perceiver \& Perceiver IO}
\fancyhead[R]{\small Dispensa da presentazione}
\fancyfoot[C]{\small\thepage}
\renewcommand{\headrulewidth}{0.4pt}

\setcounter{tocdepth}{2}
\setcounter{secnumdepth}{3}
\setlength{\parindent}{0pt}
\setlength{\parskip}{6pt plus 2pt minus 1pt}
\setlist[itemize]{itemsep=2pt,topsep=4pt,leftmargin=1.4em}
\setlist[enumerate]{itemsep=2pt,topsep=4pt,leftmargin=1.7em}

\newcommand{\R}{\mathbb{R}}
\newcommand{\bigO}{\mathcal{O}}
\newcommand{\softmax}{\operatorname{softmax}}
\newcommand{\slide}[1]{\textcolor{mainblue}{\textbf{Slide #1}}}

\newtcolorbox{takeawaybox}[1][]{
  enhanced,breakable,
  colback=softgreen,colframe=accentgreen,
  boxrule=0.8pt,arc=2pt,
  left=7pt,right=7pt,top=6pt,bottom=6pt,
  fonttitle=\bfseries,
  title={Da ricordare},
  #1
}

\newtcolorbox{formulaBox}[1][]{
  enhanced,breakable,
  colback=softblue,colframe=mainblue,
  boxrule=0.7pt,arc=2pt,
  left=7pt,right=7pt,top=6pt,bottom=6pt,
  fonttitle=\bfseries,
  title={Formula},
  #1
}

\newtcolorbox{examBox}[1][]{
  enhanced,breakable,
  colback=softamber,colframe=accentamber,
  boxrule=0.8pt,arc=2pt,
  left=7pt,right=7pt,top=6pt,bottom=6pt,
  fonttitle=\bfseries,
  title={Per l'esame},
  #1
}

\newtcolorbox{riskBox}[1][]{
  enhanced,breakable,
  colback=softred,colframe=accentred,
  boxrule=0.8pt,arc=2pt,
  left=7pt,right=7pt,top=6pt,bottom=6pt,
  fonttitle=\bfseries,
  title={Criticit\`a},
  #1
}

\title{\textbf{\textcolor{deepblue}{Perceiver \& Perceiver IO}}\\
  \large Dispensa in ordine presentazione}
\author{Francesco Gigli e Corso Vignoli}
\date{}

\begin{document}
\maketitle

\begin{center}
\begin{minipage}{0.88\linewidth}
\small
Questa dispensa ricostruisce in forma testuale la presentazione su
\textbf{Perceiver} e \textbf{Perceiver IO}. L'ordine segue la mappa del file
\texttt{Note presentazione.pptx}; le informazioni sono integrate dai sorgenti
LaTeX/PDF delle slide e dai due paper principali: Jaegle et al. (2021) e
Jaegle et al. (2022).
\end{minipage}
\end{center}

\vspace{0.5em}
\begin{takeawaybox}[title={Idea in 60 secondi}]
Il Perceiver nasce per risolvere una tensione classica: i Transformer sono
generali ma costano quadraticamente nella lunghezza dell'input, mentre le CNN
sono scalabili ma incorporano forti assunzioni sul dominio. Il Perceiver
introduce un \emph{latent array} di dimensione fissa: i latenti leggono l'input
tramite cross-attention, poi vengono raffinati da self-attention nello spazio
latente. Perceiver IO aggiunge un decoder a \emph{output queries}, cosi' la
stessa architettura puo' produrre anche output strutturati.
\end{takeawaybox}

{\small\tableofcontents}
\clearpage

% =====================================================================
\section{Introduzione e motivazione}
% =====================================================================

\subsection{Il problema: architetture specifiche per modalit\`a}
\slide{3}

Lo stato dell'arte tradizionale tende ad associare a ogni modalit\`a una
famiglia architetturale diversa:
\begin{itemize}
  \item immagini $\rightarrow$ CNN come ResNet o EfficientNet;
  \item testo $\rightarrow$ Transformer come BERT o GPT;
  \item audio $\rightarrow$ CNN 1D, RNN o modelli ibridi;
  \item point cloud $\rightarrow$ architetture specializzate come PointNet o DGCNN.
\end{itemize}

Queste scelte funzionano bene perche' incorporano \textbf{bias induttivi}
utili: localit\`a spaziale nelle immagini, ordine sequenziale nel testo,
struttura geometrica nei point cloud. Il prezzo e' che ogni nuova modalit\`a
richiede di ridisegnare parti del modello. Il Perceiver parte dalla domanda
opposta: quanto possiamo ridurre le assunzioni architetturali senza perdere
scalabilit\`a?

Il Transformer sarebbe un candidato naturale per un modello generale, perche'
l'attention non assume una griglia o una topologia specifica. Il problema e'
il costo della self-attention. Con un input di lunghezza $N$, la matrice di
attenzione ha dimensione $N \times N$:
\[
  \text{SelfAttn}(X) = \softmax\!\left(\frac{QK^\top}{\sqrt{d_k}}\right)V,
  \qquad Q,K,V \in \R^{N \times d_k}.
\]
Il costo cresce come $\bigO(N^2 d)$. Un'immagine $224 \times 224$ trattata a
livello di pixel ha $N=50\,176$: l'attenzione completa produrrebbe oltre
$2.5$ miliardi di interazioni.

\begin{riskBox}[title={Il collo di bottiglia dei Transformer}]
La generalit\`a del Transformer non basta se l'input e' molto grande. Senza un
meccanismo di compressione o sparsificazione, la self-attention piena diventa
impraticabile per immagini ad alta risoluzione, video, audio lungo o point cloud
densi.
\end{riskBox}

\subsection{La soluzione Perceiver}
\slide{4}

Il Perceiver \cite{jaegle2021perceiver} sostituisce la self-attention diretta
sull'input con un meccanismo asimmetrico. Invece di far interagire tutti gli
elementi dell'input tra loro, introduce un insieme compatto di vettori appresi,
il \textbf{latent array}. Se l'input ha $N$ elementi e i latenti sono $M$, con
$M \ll N$, i latenti leggono l'input tramite cross-attention:
\[
  Q = X_{\text{latent}} W_Q \in \R^{M \times d_k}, \qquad
  K = X_{\text{input}} W_K \in \R^{N \times d_k}, \qquad
  V = X_{\text{input}} W_V \in \R^{N \times d_v}.
\]
La matrice di attenzione e' rettangolare, $M \times N$, non piu' $N \times N$.
Il costo principale passa da $\bigO(N^2)$ a $\bigO(NM)$.

Dopo la lettura dell'input, la computazione profonda avviene solo nello spazio
latente: una pila di blocchi Transformer elabora $M$ latenti, quindi costa
$\bigO(M^2 L)$, dove $L$ e' il numero di layer latenti. Poiche' $M$ e' fissato
dall'architettura, la profondit\`a viene disaccoppiata dalla dimensione
dell'input.

\begin{formulaBox}[title={Complessit\`a del Perceiver}]
\[
  \underbrace{\bigO(N^2 d L)}_{\text{Transformer standard}}
  \quad \longrightarrow \quad
  \underbrace{\bigO(NM d)}_{\text{cross-attention input--latenti}}
  +
  \underbrace{\bigO(M^2 d L)}_{\text{self-attention latente}}.
\]
Con $M \ll N$, il termine dipendente dall'input diventa lineare in $N$.
\end{formulaBox}

Nel progetto la stessa idea viene implementata in PyTorch e testata su tre
modalit\`a: immagini, point cloud e testo. Questo e' importante perche' non si
sta solo ripetendo la teoria del paper: si verifica anche quanto l'architettura
regga in condizioni di calcolo molto piu' limitate rispetto agli esperimenti
originali di DeepMind.

\subsection{Perceiver IO: output strutturati}
\slide{5}

Il Perceiver originale e' gia' generale rispetto agli input, ma il suo output
principale e' una rappresentazione aggregata, adatta soprattutto a compiti di
classificazione. Perceiver IO \cite{jaegle2022perceiverio} estende il modello
anche dal lato degli output. L'architettura assume una forma
\textbf{read-process-write}:
\begin{enumerate}
  \item \textbf{Read}: l'input viene letto dai latenti tramite cross-attention.
  \item \textbf{Process}: i latenti vengono raffinati da self-attention profonda.
  \item \textbf{Write}: un insieme di \emph{output queries} interroga i latenti
  e produce l'output desiderato.
\end{enumerate}

La differenza concettuale e' nel decoder. Ogni output viene associato a una
query che ne specifica la semantica. Per classificare basta una singola query;
per un output denso, come optical flow, si possono usare query indicizzate dalle
coordinate dei pixel; per masked language modeling si usa una query per ogni
posizione mascherata.

\begin{takeawaybox}
Perceiver IO non cambia il cuore del Perceiver: cambia il modo in cui si legge
lo spazio latente in uscita. L'encoder resta agnostico alla modalit\`a, mentre
le output queries rendono flessibile la forma dell'output.
\end{takeawaybox}

% =====================================================================
\section{Architettura}
% =====================================================================

\subsection{Panoramica completa}
\slide{7}

L'architettura puo' essere letta come una pipeline con quattro oggetti:
\begin{itemize}
  \item \textbf{Input array} $X \in \R^{N \times C}$: sequenza generica ottenuta
  da pixel, punti 3D, byte di testo o altri elementi;
  \item \textbf{positional encoding}: coordinate o Fourier features concatenate
  ai canali originali;
  \item \textbf{latent array} $Z \in \R^{M \times D}$: bottleneck appreso,
  compatto e indipendente dalla lunghezza dell'input;
  \item \textbf{output head/decoder}: pooling nel Perceiver base, output
  queries nel Perceiver IO.
\end{itemize}

\begin{center}
\begin{tikzpicture}[
  node distance=0.8cm,
  block/.style={draw,rounded corners=2pt,thick,align=center,
    minimum width=2.25cm,minimum height=0.85cm,font=\small},
  arr/.style={-Stealth,thick}
]
\node[block,fill=softblue] (input) {Input\\$N \times C$};
\node[block,fill=softamber,right=of input] (pe) {Fourier PE\\$N \times C_{\text{tot}}$};
\node[block,fill=softgreen,right=of pe] (cross) {Cross-attn\\$M \times N$};
\node[block,fill=softblue,right=of cross] (latent) {Latent\\Transformer};
\node[block,fill=softgreen,right=of latent] (out) {Output\\head/queries};
\draw[arr] (input) -- (pe);
\draw[arr] (pe) -- (cross);
\draw[arr] (cross) -- (latent);
\draw[arr] (latent) -- (out);
\node[below=0.35cm of cross,font=\scriptsize,align=center,text=deepblue]
  {i latenti interrogano l'input};
\end{tikzpicture}
\end{center}

La propriet\`a da ricordare e' che l'input puo' cambiare lunghezza e modalit\`a,
ma la parte profonda del modello opera sempre su $M$ latenti. Questo rende
l'architettura adatta a input grandi e multimodali, pur usando blocchi
Transformer standard nello spazio compresso.

\subsection{Cross-attention: il meccanismo chiave}
\slide{8}

Nella self-attention standard, query, key e value provengono dallo stesso
array. Nella cross-attention del Perceiver, invece, le query provengono dai
latenti, mentre key e value provengono dall'input:
\[
  Q = ZW_Q, \qquad K = XW_K, \qquad V = XW_V.
\]
Il risultato e' una lettura selettiva dell'input:
\[
  \operatorname{CrossAttn}(Z,X)
  =
  \softmax\!\left(\frac{QK^\top}{\sqrt{d_k}}\right)V.
\]

Ogni riga della matrice $QK^\top$ corrisponde a un latente; ogni colonna a un
elemento dell'input. Il latente decide quali parti dell'input leggere, aggrega
i value con pesi di attenzione e aggiorna la propria rappresentazione. L'input
non deve interagire direttamente con tutti gli altri input: l'informazione passa
attraverso il bottleneck latente.

\begin{examBox}
Se ti chiedono ``perche' la complessit\`a diventa lineare?'', la risposta breve
e': perche' la matrice di attenzione non e' piu' $N \times N$, ma $M \times N$.
Dato che $M$ e' fissato e molto piu' piccolo di $N$, il costo cresce
linearmente con la lunghezza dell'input.
\end{examBox}

\subsection{Latent Transformer e weight sharing}
\slide{9}

Dopo la cross-attention, i latenti contengono informazione letta dall'input.
Serve pero' una fase di elaborazione interna: i latenti devono comunicare tra
loro, comporre feature e raffinare la rappresentazione. Questa fase e' il
\textbf{latent Transformer}: una sequenza di blocchi pre-norm con
self-attention, residual connections e MLP espansiva.

Poiche' la self-attention avviene solo sui latenti, il costo e'
$\bigO(M^2)$, non $\bigO(N^2)$. Nel progetto, ad esempio, una configurazione
CIFAR-10 usa $M=96$ latenti e profondit\`a ricorrente, una dimensione molto
piu' gestibile dell'input completo.

Il \textbf{weight sharing} riapplica lo stesso blocco per piu' iterazioni. La
lettura intuitiva e' simile a un RNN nello spazio latente: lo stesso processore
raffina lo stato piu' volte. Il vantaggio e' una forte riduzione dei parametri;
lo svantaggio e' una possibile perdita di capacit\`a rispetto a blocchi tutti
indipendenti.

\begin{center}
\begin{tabular}{lcc}
\toprule
\textbf{Configurazione CIFAR-10} & \textbf{Parametri} & \textbf{Accuratezza} \\
\midrule
Shared baseline & $3.35$M & $69.69\%$ \\
Weight-sharing control & $3.35$M & $68.70\%$ \\
No sharing & $11.0$M & $73.61\%$ \\
\bottomrule
\end{tabular}
\end{center}

\begin{takeawaybox}
Il weight sharing e' un compromesso: meno parametri e memoria, ma meno
capacit\`a rispetto a layer separati. Nel progetto produce circa un terzo dei
parametri della versione non condivisa, mantenendo prestazioni competitive.
\end{takeawaybox}

\subsection{Fourier positional encoding}
\slide{10}

L'attention e' permutation-equivariant: se permutiamo gli elementi di input e
non forniamo informazioni di posizione, il modello non sa dove si trovino pixel,
punti o token. Per immagini e point cloud questa informazione e' fondamentale.
Il Perceiver concatena ai canali originali un positional encoding basato su
Fourier features.

Per una coordinata $x_i$ e una frequenza $f_k$, si usano termini sinusoidali:
\[
  \gamma(x_i)
  =
  \left[
    \sin(2\pi f_k x_i),\;
    \cos(2\pi f_k x_i)
  \right]_{k=1}^{K}.
\]
Nel documento delle slide, per immagini si ottengono $2$ coordinate grezze piu'
$64$ componenti Fourier, concatenate ai canali RGB; per point cloud si usano
coordinate 3D; per testo una coordinata posizionale lungo la sequenza.

\begin{center}
\begin{tabular}{lcc}
\toprule
\textbf{Setting CIFAR-10} & \textbf{Accuratezza} & \textbf{Differenza} \\
\midrule
Con Fourier PE & $72.23\%$ & baseline \\
Solo RGB, senza PE & $35.41\%$ & $-36.82$ pp \\
\bottomrule
\end{tabular}
\end{center}

\begin{takeawaybox}[title={Risultato sperimentale piu' forte}]
Il positional encoding e' il componente piu' critico nella riproduzione su
CIFAR-10: senza posizione il modello tratta l'immagine quasi come un insieme
non ordinato di patch/pixel e l'accuratezza si dimezza.
\end{takeawaybox}

\subsection{Decoder Perceiver IO con output queries}
\slide{11}

Nel Perceiver IO il decoder usa una cross-attention opposta rispetto
all'encoder. Nell'encoder, i latenti interrogano l'input; nel decoder, le
output queries interrogano i latenti:
\[
  Q_{\text{dec}} = Q_{\text{out}} W_Q, \qquad
  K_{\text{dec}} = Z W_K, \qquad
  V_{\text{dec}} = Z W_V.
\]
L'output e':
\[
  Y =
  \softmax\!\left(\frac{Q_{\text{dec}}K_{\text{dec}}^\top}{\sqrt{d_k}}\right)
  V_{\text{dec}}
  \in \R^{O \times E},
\]
dove $O$ e' il numero di query di output. Questo numero non deve coincidere ne'
con la lunghezza dell'input $N$ ne' con il numero di latenti $M$.

Questa separazione e' il punto di forza di Perceiver IO: la rete puo' mantenere
uno spazio latente compatto, ma produrre output di forma molto diversa a seconda
del compito.

\subsection{Design delle output queries per task}
\slide{12}

\begin{center}
\begin{tabularx}{\linewidth}{l l X}
\toprule
\textbf{Task} & \textbf{Query} & \textbf{Output} \\
\midrule
Classificazione & $O=1$, vettore appreso & logit di classe \\
Optical flow & $O=H \times W$, coordinate dei pixel & vettori di flusso per pixel \\
Masked LM & una query per posizione mascherata & distribuzione sul vocabolario \\
Multimodale & query/task embedding variabili & output misti o multi-task \\
\bottomrule
\end{tabularx}
\end{center}

Nel progetto, la classificazione su CIFAR-10 e ModelNet40 usa una singola query
o una testa equivalente di aggregazione; il masked language modeling su
WikiText-103 usa query associate alle posizioni mascherate e predice byte su un
vocabolario di dimensione $256$.

\subsection{Analisi della complessit\`a computazionale}
\slide{13}

La complessit\`a dei tre casi principali puo' essere riassunta cosi':
\[
\begin{aligned}
  \text{Transformer standard} &: \bigO(N^2 d L),\\
  \text{Perceiver} &: \bigO((NM + M^2L)d),\\
  \text{Perceiver IO} &: \bigO(((N+O)M + M^2L)d).
\end{aligned}
\]
Il termine aggiuntivo di Perceiver IO e' $OM$, dovuto al decoder che fa
attendere le output queries ai latenti.

Esempio dalle slide: con $N=1024$, $M=96$, $L=4$,
\[
  N^2 = 1\,048\,576,
  \qquad
  NM + M^2L = 1024 \cdot 96 + 96^2 \cdot 4 = 135\,168.
\]
La riduzione e' circa $7.8\times$. Su immagini piu' grandi il vantaggio cresce:
con $224 \times 224$ elementi, un Transformer pieno arriva a circa
$2.5 \times 10^9$ interazioni, mentre un Perceiver con $M=512$ resta nell'ordine
di decine di milioni.

\subsection{GELU e ottimizzatore LAMB}
\slide{14}

Nei blocchi Transformer si usa una MLP con attivazione GELU:
\[
  \operatorname{GELU}(x) = x \Phi(x)
  \approx x \sigma(1.702x).
\]
GELU e' liscia e non monotona, ed e' diventata standard in molte architetture
Transformer. Alcune varianti usano anche GEGLU, cioe' una versione gated:
\[
  \operatorname{GEGLU}(x)
  =
  (xW_1) \odot \operatorname{GELU}(xW_2).
\]

Per l'ottimizzazione il progetto usa LAMB \cite{you2020lamb}, un ottimizzatore
pensato per training a batch grandi. L'idea e' scalare l'update layer-wise con
un \emph{trust ratio}, dipendente dalla norma dei pesi e dalla norma
dell'update. Questo aiuta quando gradienti e scale dei layer sono molto diversi.

\begin{riskBox}
LAMB funziona al meglio con batch grandi. Nel progetto, i vincoli hardware
limitano batch size e numero di latenti; questo spiega parte del divario tra i
risultati originali e la riproduzione.
\end{riskBox}

\subsection{Implementazione: struttura del codice}
\slide{15}

La presentazione descrive un'implementazione \emph{from scratch} in PyTorch e
PyTorch Lightning. La divisione logica e':
\begin{itemize}
  \item modelli: Perceiver, Perceiver IO, attention, positional encoding,
  ottimizzatore LAMB;
  \item data module: CIFAR-10, ModelNet40, WikiText-103, GLUE;
  \item script esperimenti: training e ablation per immagini, point cloud,
  testo;
  \item utility: analisi, riproducibilit\`a e generazione risultati.
\end{itemize}

La dimensione indicativa e' circa $2500$ linee: abbastanza piccola da restare
leggibile, ma completa per verificare tre modalit\`a e piu' ablation.

% =====================================================================
\section{Esperimenti e risultati}
% =====================================================================

\subsection{CIFAR-10: setup sperimentale}
\slide{17}

L'obiettivo su CIFAR-10 e' validare il Perceiver su immagini senza usare bias
convoluzionali. Il modello riceve patch/pixel con positional encoding e deve
imparare da dati relativamente piccoli rispetto a ImageNet.

\begin{center}
\begin{tabular}{ll}
\toprule
\textbf{Parametro} & \textbf{Valore} \\
\midrule
Latenti & $96 \times 384$ \\
Cross-attention stages & $4$ \\
Transformer blocks & $4$, condivisi \\
Attention heads & $3$ \\
Patch size & $2 \times 2$ \\
Dropout & $0.2$ \\
Optimizer & LAMB \\
Learning rate & $0.004$ \\
Batch size & $64$ \\
Epochs & $120$ \\
Parametri & circa $3.35$M \\
\bottomrule
\end{tabular}
\end{center}

Le ablation principali sono: rimozione o variazione del positional encoding,
weight sharing contro blocchi non condivisi, robustezza alla permutazione dei
pixel e confronto Perceiver base contro Perceiver IO.

\subsection{CIFAR-10: risultati delle ablation}
\slide{18}

\begin{center}
\begin{tabular}{lc}
\toprule
\textbf{Configurazione} & \textbf{Accuratezza} \\
\midrule
Perceiver IO & $78.20\%$ \\
No weight sharing & $73.61\%$ \\
Fourier control & $72.23\%$ \\
Baseline Fourier & $69.69\%$ \\
Weight-sharing control & $68.70\%$ \\
Fourier permuted & $62.31\%$ \\
Learned PE permuted & $55.12\%$ \\
RGB only, senza PE & $35.41\%$ \\
\bottomrule
\end{tabular}
\end{center}

Tre osservazioni emergono chiaramente:
\begin{enumerate}
  \item il positional encoding e' indispensabile;
  \item Perceiver IO migliora il Perceiver base grazie a una lettura piu'
  flessibile dello spazio latente;
  \item il weight sharing riduce fortemente i parametri, ma perde qualche punto
  rispetto a blocchi separati.
\end{enumerate}

\subsection{Impatto del positional encoding}
\slide{19}

Senza positional encoding il modello non distingue la struttura spaziale. In
termini pratici, un'immagine diventa un insieme di valori RGB: l'informazione su
``dove'' si trova ogni colore viene persa. Questo spiega il crollo da
$72.23\%$ a $35.41\%$.

Il risultato conferma l'ipotesi del paper originale: un'architettura molto
generale puo' evitare bias forti come convoluzioni o pooling locali, ma non puo'
rinunciare completamente alla posizione. La posizione non e' un dettaglio: e'
parte del dato.

\subsection{Weight sharing: efficienza contro accuratezza}
\slide{20}

Il confronto piu' diretto e':
\begin{itemize}
  \item shared baseline: $3.35$M parametri, $69.69\%$;
  \item no sharing: $11.0$M parametri, $73.61\%$.
\end{itemize}
La versione non condivisa guadagna $3.92$ punti percentuali, ma usa circa
$3.3\times$ piu' parametri. In termini di efficienza, il weight sharing
mantiene una quota significativa della performance usando molte meno risorse.

\begin{examBox}
Una buona risposta orale: ``Il weight sharing trasforma il blocco in un
processore ricorrente dei latenti. Riduce i parametri e regolarizza, ma limita
la capacit\`a perche' ogni iterazione deve usare gli stessi pesi.''
\end{examBox}

\subsection{Robustezza alla permutazione dei pixel}
\slide{21}

L'esperimento applica una permutazione fissa alle posizioni dei pixel. I
risultati mostrano che Fourier PE degrada meno del positional encoding appreso:
\begin{center}
\begin{tabular}{lcc}
\toprule
\textbf{Configurazione} & \textbf{Accuratezza} & \textbf{Drop} \\
\midrule
Fourier PE normale & $72.23\%$ & -- \\
Fourier PE permutato & $62.31\%$ & $-9.92$ pp \\
Learned PE permutato & $55.12\%$ & $-17.11$ pp \\
\bottomrule
\end{tabular}
\end{center}

L'interpretazione e' che le Fourier features codificano regolarit\`a
multi-scala e mantengono una struttura piu' robusta, mentre un positional
encoding appreso tende a memorizzare posizioni assolute e soffre di piu' quando
la geometria viene alterata.

\subsection{Perceiver contro Perceiver IO su CIFAR-10}
\slide{22}

Su CIFAR-10, Perceiver IO raggiunge $78.20\%$, contro $72.23\%$ del Perceiver
con Fourier PE. Il guadagno e' di $5.97$ punti percentuali.

La ragione e' che l'output query agisce come aggregatore appreso. Il Perceiver
base usa una forma piu' semplice di pooling sui latenti; Perceiver IO, invece,
lascia che una query impari quali latenti leggere e come combinarli per il
task. Questo aggiunge capacit\`a e rende la lettura dello spazio latente piu'
flessibile.

\subsection{ModelNet40: setup e risultati}
\slide{23}

ModelNet40 testa la capacit\`a del modello su point cloud 3D. Qui il Perceiver
e' particolarmente naturale: un point cloud e' un insieme di punti, non una
griglia regolare come un'immagine.

\begin{center}
\begin{tabular}{ll}
\toprule
\textbf{Parametro} & \textbf{Valore} \\
\midrule
Latenti & $128 \times 512$ \\
Cross-attention stages & $2$ \\
Transformer blocks & $6$, condivisi \\
Punti per oggetto & $2048$ \\
Optimizer & LAMB \\
Epochs & $200$ \\
Batch size & $128$ \\
\bottomrule
\end{tabular}
\end{center}

Risultati:
\begin{center}
\begin{tabular}{lc}
\toprule
\textbf{Configurazione} & \textbf{Accuratezza} \\
\midrule
Baseline con scaling & $84.16\%$ \\
Scaling + rotation & $83.06\%$ \\
Scaling + translation & $82.90\%$ \\
Paper originale & $85.7\%$ \\
\bottomrule
\end{tabular}
\end{center}

Il gap rispetto al paper e' solo $1.54$ punti percentuali, un risultato molto
buono considerando batch e hardware piu' limitati.

\subsection{Effetti della data augmentation su ModelNet40}
\slide{24}

In questo caso, l'augmentation aggressiva peggiora leggermente. La spiegazione
plausibile e' tripla:
\begin{itemize}
  \item il modello ha capacit\`a limitata rispetto alla configurazione originale;
  \item rotazioni e traslazioni richiedono piu' dati o piu' epoche per essere
  assorbite;
  \item il positional encoding Fourier usa coordinate assolute, quindi una
  trasformazione geometrica sposta anche la codifica posizionale.
\end{itemize}

Il punto non e' che l'augmentation sia sempre negativa, ma che va calibrata con
capacita' del modello, budget di training e tipo di positional encoding.

\subsection{WikiText-103: pre-training MLM byte-level}
\slide{25}

Per il testo, la scelta piu' coerente con l'idea di architettura
modality-agnostic e' usare input byte-level. Il vocabolario ha dimensione
$256$, una per ogni byte, e non richiede tokenizzatori BPE o WordPiece.

\begin{center}
\begin{tabular}{ll}
\toprule
\textbf{Aspetto} & \textbf{Valore} \\
\midrule
Dataset & WikiText-103 \\
Tokenizzazione & byte-level UTF-8 \\
Vocabolario & $256$ \\
Sequence length & $1024$ \\
Mask probability & $15\%$ \\
Modello & Perceiver IO \\
Parametri & circa $10$M \\
Epochs & $50$ \\
\bottomrule
\end{tabular}
\end{center}

Il vantaggio e' la generalit\`a: niente tokenizzatore specifico per lingua o
dominio. Lo svantaggio e' che le sequenze diventano piu' lunghe e il modello
non riceve direttamente unit\`a subword semanticamente dense.

\subsection{GLUE: fine-tuning su 8 task}
\slide{26}

La pipeline di transfer learning e':
\begin{enumerate}
  \item pre-training MLM su WikiText-103;
  \item trasferimento dei pesi dell'encoder;
  \item sostituzione delle output queries con una query di classificazione;
  \item fine-tuning sui task GLUE.
\end{enumerate}

I task coprono fenomeni diversi:
\begin{itemize}
  \item singola frase: CoLA, SST-2;
  \item similarit\`a/parafrasi: MRPC, STS-B, QQP;
  \item inferenza linguistica: MNLI, QNLI, RTE.
\end{itemize}

Il setup di fine-tuning indicato dalle slide usa $30$ epoche, learning rate
$5 \times 10^{-4}$ e input byte-level.

\subsection{Analisi per task GLUE}
\slide{27}

I risultati sono sopra la majority baseline sulla maggior parte dei task, con
buona performance su SST-2 e maggiori difficolt\`a su task piu' sottili come
CoLA o MNLI. Questo e' coerente con le limitazioni del modello:
\begin{itemize}
  \item il byte-level rende le sequenze piu' lunghe;
  \item il modello e' molto piu' piccolo di BERT base;
  \item manca il vantaggio semantico di una tokenizzazione subword;
  \item il pre-training e' limitato dal budget computazionale.
\end{itemize}

Il messaggio non e' che Perceiver IO batta sempre BERT in questa riproduzione,
ma che la stessa architettura puo' passare da immagini e point cloud al testo
senza cambiare il meccanismo centrale.

\subsection{Convergenza tra modalit\`a}
\slide{28}

Le tre modalit\`a convergono con tempi diversi:
\begin{center}
\begin{tabular}{lc}
\toprule
\textbf{Modalit\`a} & \textbf{Epoche indicative} \\
\midrule
Point cloud & circa $50$ \\
GLUE fine-tuning & $15$--$30$ \\
CIFAR-10 & circa $120$ \\
\bottomrule
\end{tabular}
\end{center}

I point cloud hanno input relativamente basso-dimensionali e convergono piu'
rapidamente. GLUE beneficia del pre-training. CIFAR-10, invece, e' difficile
perche' il modello vede informazione molto grezza e non usa bias convoluzionali.

\subsection{Attention maps}
\slide{29}

Le attention maps mostrano qualitativamente l'effetto del positional encoding.
Con Fourier PE, la cross-attention impara pattern spazialmente coerenti; senza
PE, l'entropia resta piu' alta e i pattern appaiono meno strutturati. Questo
supporta la lettura quantitativa: la posizione permette ai latenti di
specializzarsi su regioni o strutture dell'input.

\subsection{Riassunto complessivo degli esperimenti}
\slide{30}

\begin{center}
\begin{tabularx}{\linewidth}{l X l}
\toprule
\textbf{Modalit\`a} & \textbf{Risultato principale} & \textbf{Nota} \\
\midrule
Immagini & Perceiver $72.23\%$, Perceiver IO $78.20\%$ & PE decisivo \\
Point cloud & ModelNet40 $84.16\%$ & vicino al paper \\
Testo & MLM + GLUE fine-tuning & byte-level, 8 task \\
\bottomrule
\end{tabularx}
\end{center}

Nel complesso: oltre 20 training run, tre modalit\`a, ablation su componenti
chiave e una pipeline riproducibile. Il valore della riproduzione e' mostrare
non solo cosa funziona nel paper, ma cosa resta solido con risorse limitate.

% =====================================================================
\section{Criticit\`a e discussione}
% =====================================================================

\subsection{Gap rispetto al paper originale}
\slide{32}

Il gap e' molto diverso tra point cloud e immagini:
\begin{center}
\begin{tabular}{lccc}
\toprule
\textbf{Dataset} & \textbf{Paper} & \textbf{Progetto} & \textbf{Gap} \\
\midrule
ModelNet40 & $85.7\%$ & $84.16\%$ & $1.54$ pp \\
CIFAR-10 / immagini & $>90\%$ & $78.20\%$ & circa $12$ pp \\
\bottomrule
\end{tabular}
\end{center}

Su ModelNet40 la riproduzione e' molto vicina. Su immagini il gap e'
significativo e dipende soprattutto dal budget computazionale: batch piu'
piccolo, meno latenti, modello piu' piccolo e meno tempo di training.

\subsection{Limitazioni hardware e scelte implementative}
\slide{33}

Il confronto con gli esperimenti originali e' sbilanciato:
\begin{center}
\begin{tabularx}{\linewidth}{l X X}
\toprule
\textbf{Aspetto} & \textbf{Paper DeepMind} & \textbf{Progetto} \\
\midrule
Hardware & centinaia di TPU v3 & una GPU consumer \\
Batch size & fino a $1024$ & $64$--$128$ \\
Latenti & $256$--$512$ & $96$--$128$ \\
Parametri & decine/centinaia di milioni & $3.35$--$10$M \\
\bottomrule
\end{tabularx}
\end{center}

Le conseguenze sono dirette: LAMB e' meno efficace con batch piccoli, meno
latenti riducono la capacit\`a del bottleneck, e i modelli piu' piccoli faticano
di piu' sulle immagini grezze.

\subsection{Possibili miglioramenti}
\slide{34}

Le direzioni di miglioramento piu' naturali sono:
\begin{itemize}
  \item aumentare il numero di latenti con gradient checkpointing;
  \item usare cross-attention multi-scala, da rappresentazioni coarse a fine;
  \item esplorare positional encoding relativi, come RoPE o ALiBi;
  \item aumentare il batch effettivo con gradient accumulation;
  \item usare mixed precision FP16/BF16;
  \item introdurre augmentation moderne su immagini, come CutMix e MixUp;
  \item usare tokenizzazione subword per confronti NLP piu' forti;
  \item confrontare con efficient Transformer come Linformer o Performer.
\end{itemize}

\subsection{Lezioni apprese}
\slide{35}

Le lezioni principali sono quattro:
\begin{enumerate}
  \item il Perceiver e' realmente modality-agnostic nella struttura;
  \item Perceiver IO migliora la lettura dello spazio latente;
  \item il positional encoding e' il componente piu' critico;
  \item il compute disponibile determina fortemente il gap sulle immagini.
\end{enumerate}

\begin{takeawaybox}[title={Takeaway principale}]
La generalit\`a non elimina la necessit\`a di informazione strutturale. Il
Perceiver evita convoluzioni, RNN e moduli specifici, ma deve comunque sapere
dove si trovano gli elementi dell'input.
\end{takeawaybox}

\subsection{Conclusioni}
\slide{36}

Il Perceiver dimostra che una singola architettura attentiva puo' trattare
modalit\`a diverse mantenendo costo lineare nella dimensione dell'input. Il
latent bottleneck rende trattabili input grandi, mentre il latent Transformer
mantiene la potenza espressiva dei Transformer.

Perceiver IO completa il quadro rendendo generali anche gli output. Con le
output queries, lo stesso spazio latente puo' essere letto per classificazione,
language modeling, output densi o task multimodali.

I limiti restano importanti: il modello richiede molto compute per competere
con architetture specializzate su immagini; dipende fortemente dal positional
encoding; e il bottleneck latente puo' perdere dettagli fini se troppo piccolo.
Tuttavia, come direzione di ricerca, e' un passo convincente verso modelli di
percezione general-purpose.

% =====================================================================
\section{Riferimenti}
% =====================================================================
\slide{37}

\begin{thebibliography}{9}
\bibitem{jaegle2021perceiver}
Andrew Jaegle, Felix Gimeno, Andrew Brock, Andrew Zisserman, Oriol Vinyals,
Joao Carreira.
\newblock \emph{Perceiver: General Perception with Iterative Attention}.
\newblock ICML, 2021.

\bibitem{jaegle2022perceiverio}
Andrew Jaegle et al.
\newblock \emph{Perceiver IO: A General Architecture for Structured Inputs and Outputs}.
\newblock ICLR, 2022.

\bibitem{vaswani2017attention}
Ashish Vaswani et al.
\newblock \emph{Attention Is All You Need}.
\newblock NeurIPS, 2017.

\bibitem{you2020lamb}
Yang You et al.
\newblock \emph{Large Batch Optimization for Deep Learning: Training BERT in 76 Minutes}.
\newblock ICLR, 2020.

\bibitem{tancik2020fourier}
Matthew Tancik et al.
\newblock \emph{Fourier Features Let Networks Learn High Frequency Functions}.
\newblock NeurIPS, 2020.

\bibitem{shazeer2020glu}
Noam Shazeer.
\newblock \emph{GLU Variants Improve Transformer}.
\newblock arXiv, 2020.
\end{thebibliography}

% =====================================================================
\appendix
\section{Domande attese all'esame}
\slide{38}

\subsection{Perche' cross-attention invece di self-attention piena?}
Per ridurre la complessit\`a da $\bigO(N^2)$ a $\bigO(NM)$. I latenti sono
molto meno numerosi degli elementi di input, quindi l'attenzione rettangolare
$M \times N$ e' molto piu' economica della matrice quadrata $N \times N$.

\subsection{Che effetto ha il weight sharing?}
Riduce i parametri e rende il blocco simile a un processore ricorrente nello
spazio latente. Nel progetto produce circa $3\times$ meno parametri rispetto
alla versione non condivisa, ma perde alcuni punti di accuratezza.

\subsection{Perche' Fourier PE e non solo learned PE?}
Le Fourier features danno una codifica multi-scala delle coordinate e risultano
piu' robuste quando la struttura spaziale viene perturbata. Il learned PE puo'
funzionare, ma tende a legarsi alle posizioni assolute viste in training.

\subsection{Perche' il gap su CIFAR-10 e' grande?}
Perche' la configurazione riprodotta e' molto piu' piccola di quella originale:
batch size minore, meno latenti, meno parametri e molto meno compute. Inoltre
le immagini grezze sono difficili senza bias convoluzionali.

\subsection{Perche' usare byte-level tokenization?}
Per restare coerenti con l'idea modality-agnostic: un vocabolario di $256$ byte
non richiede tokenizzatori specifici. Il costo e' avere sequenze piu' lunghe e
meno informazione semantica pre-strutturata rispetto a BPE/WordPiece.

\subsection{Come migliora Perceiver IO rispetto al Perceiver base?}
Perceiver IO sostituisce un readout semplice con output queries apprese o
strutturate. Questo permette di interrogare i latenti in modo task-specific e
di produrre output con forma arbitraria.

\subsection{Puo' sostituire ViT o BERT?}
Non necessariamente nelle configurazioni piccole. Il valore del Perceiver e'
la generalit\`a e la scalabilit\`a su input grandi e multimodali. Per competere
con modelli specializzati servono compute, dati e dimensioni comparabili.

\section*{Chiusura}
\slide{39}

\begin{center}
\Large\textbf{Perceiver = input grandi, architettura unica, bottleneck latente.}\\[0.5em]
\large Perceiver IO = la stessa idea, ma anche per output strutturati.
\end{center}

\end{document}
